\section{Preliminaries}

\subsection{Conventions and Units}
We use natural units:
\begin{equation}
\hbar = c = 1.
\end{equation}
Mass, energy, momentum have the same dimensions:
\begin{equation}
[E]=[p]=[m]=L^{-1}.
\end{equation}

\subsection{Indices and Metric}
Spacetime coordinates:
\begin{equation}
x^\mu = (t,\mathbf{x}), \qquad \mu=0,1,2,3.
\end{equation}
Metric signature (choose one and be consistent). Common choice:
\begin{equation}
\eta_{\mu\nu}=\mathrm{diag}(1,-1,-1,-1).
\end{equation}
Scalar product:
\begin{equation}
x\cdot y := x^\mu y_\mu = x^0 y^0 - \mathbf{x}\cdot \mathbf{y}.
\end{equation}
Raising/lowering indices:
\begin{equation}
x_\mu = \eta_{\mu\nu}x^\nu.
\end{equation}

Derivatives:
\begin{equation}
\partial_\mu := \frac{\partial}{\partial x^\mu}, \qquad \partial^\mu = \eta^{\mu\nu}\partial_\nu.
\end{equation}
D'Alembertian:
\begin{equation}
\square := \partial_\mu \partial^\mu = \partial_t^2 - \nabla^2.
\end{equation}

\subsection{Lorentz Transformations}
Lorentz transformation:
\begin{equation}
x'^\mu = \Lambda^\mu_{\ \nu}x^\nu,
\end{equation}
with invariance condition:
\begin{equation}
\eta_{\rho\sigma}\Lambda^\rho_{\ \mu}\Lambda^\sigma_{\ \nu}=\eta_{\mu\nu}.
\end{equation}
Invariant interval:
\begin{equation}
x^2 = x_\mu x^\mu = t^2-|\mathbf{x}|^2.
\end{equation}

\subsection{Momentum, On-shell Condition}
Four-momentum:
\begin{equation}
p^\mu=(E,\mathbf{p}), \qquad p^2 = p_\mu p^\mu = E^2-|\mathbf{p}|^2.
\end{equation}
Mass shell condition:
\begin{equation}
p^2 = m^2 \quad \Rightarrow \quad E = \pm \sqrt{|\mathbf{p}|^2+m^2}.
\end{equation}
Usually take positive energy:
\begin{equation}
E_{\mathbf{p}}:=\sqrt{|\mathbf{p}|^2+m^2}.
\end{equation}

\subsection{Fourier Transform Conventions}
A common convention in QFT:
\begin{align}
f(x) &= \int \frac{d^4p}{(2\pi)^4}\, e^{-ip\cdot x}\,\tilde f(p),\\
\tilde f(p) &= \int d^4x\, e^{ip\cdot x}\, f(x).
\end{align}

Spatial Fourier transform:
\begin{align}
f(\mathbf{x}) &= \int \frac{d^3p}{(2\pi)^3}\, e^{i\mathbf{p}\cdot\mathbf{x}}\,\tilde f(\mathbf{p}),\\
\tilde f(\mathbf{p}) &= \int d^3x\, e^{-i\mathbf{p}\cdot\mathbf{x}}\, f(\mathbf{x}).
\end{align}

Useful identity:
\begin{equation}
\int d^4x\, e^{i(p-q)\cdot x} = (2\pi)^4\delta^{(4)}(p-q).
\end{equation}

\subsection{Delta Functions}
Dirac delta (1D):
\begin{equation}
\int_{-\infty}^{\infty} dx\, \delta(x-a) f(x) = f(a).
\end{equation}
Scaling rule:
\begin{equation}
\delta(ax)=\frac{1}{|a|}\delta(x).
\end{equation}

Multi-dimensional delta:
\begin{equation}
\delta^{(n)}(\mathbf{x})=\prod_{i=1}^n \delta(x_i),
\qquad
\int d^n x\, \delta^{(n)}(\mathbf{x}-\mathbf{a}) f(\mathbf{x})=f(\mathbf{a}).
\end{equation}

\subsection{Distributions and the $i\epsilon$ Prescription}
A key distribution identity:
\begin{equation}
\frac{1}{x\pm i\epsilon} = \mathcal{P}\left(\frac{1}{x}\right)\mp i\pi \delta(x),
\qquad \epsilon\to 0^+.
\end{equation}
Here $\mathcal{P}$ denotes the Cauchy principal value.

\subsection{Heaviside Step Function}
Definition:
\begin{equation}
\theta(t)=
\begin{cases}
1, & t>0,\\
0, & t<0.
\end{cases}
\end{equation}
Distribution derivative:
\begin{equation}
\frac{d}{dt}\theta(t)=\delta(t).
\end{equation}

\subsection{Useful Integrals and Identities}
Gaussian integral (real, $a>0$):
\begin{equation}
\int_{-\infty}^{\infty} dx\, e^{-ax^2} = \sqrt{\frac{\pi}{a}}.
\end{equation}

Integration by parts (fields vanish at infinity):
\begin{equation}
\int d^4x\, (\partial_\mu f) g = -\int d^4x\, f (\partial_\mu g).
\end{equation}

\subsection{Notation Summary}
\begin{align}
x^\mu &= (t,\mathbf{x}), &
p^\mu &= (E,\mathbf{p}),\\
x\cdot p &= x^\mu p_\mu = Et-\mathbf{p}\cdot\mathbf{x}, &
\square &= \partial_\mu\partial^\mu.
\end{align}
